\chapter{Conception de la solution}

\section*{Introduction}
Ce chapitre présente la conception retenue pour relier le dashboard à Nav2. Il décrit
l'architecture générale, justifie le choix d'une passerelle ROS~2 dédiée, puis détaille le
protocole de communication, la machine d'états d'une demande de navigation, la gestion
des repères et les règles de sûreté.

\section{Architecture générale}
La figure~\ref{fig:architecture} présente la chaîne retenue. Le navigateur communique avec
ROS~2 par rosbridge, au moyen d'une connexion WebSocket. Il n'appelle pas Nav2
directement : toutes les commandes passent par un nœud intermédiaire,
\code{dashboard\_navigation}, qui vérifie la demande, dialogue avec Nav2 par l'action
\code{/navigate\_to\_pose} et publie en retour un état synthétique. La position du robot
est lue dans l'arbre TF (transformation \code{map}~$\rightarrow$~\code{base\_link}), qui
combine la localisation AMCL et l'odométrie.

\begin{figure}[htbp]
\centering
\resizebox{\linewidth}{!}{%
\begin{tikzpicture}[node distance=0.9cm and 1.1cm]
  \node[web,minimum width=3.4cm] (ui) {Live Map (React)\\\scriptsize NavigationPanel\\\scriptsize useNavigation};
  \node[ros,right=1.6cm of ui,minimum width=2.6cm] (rb) {rosbridge\\\scriptsize WebSocket :9090};
  \node[neuf,right=1.6cm of rb,minimum width=3.4cm] (gw) {Passerelle\\\code{dashboard\_navigation}};
  \node[ros,right=1.6cm of gw,minimum width=2.6cm] (nav) {Nav2\\\scriptsize bt\_navigator};
  \node[sim,below=1.5cm of nav,minimum width=2.6cm] (gz) {Gazebo\\\scriptsize robot simulé};
  \node[ros,below=1.5cm of gw,minimum width=3.4cm] (tf) {TF \code{map}$\rightarrow$\code{base\_link}\\\scriptsize + costmap globale};
  \node[web,below=1.5cm of ui,minimum width=3.4cm,fill=gray!10] (api) {Backend FastAPI\\\scriptsize comptes, flotte};
  \draw[fl] ([yshift=2mm]ui.east) -- node[lbl,above=1pt]{services} ([yshift=2mm]rb.west);
  \draw[fl] ([yshift=-2mm]rb.west) -- node[lbl,below=1pt]{état JSON} ([yshift=-2mm]ui.east);
  \draw[fl] ([yshift=2mm]rb.east) -- node[lbl,above=1pt]{submit / cancel} ([yshift=2mm]gw.west);
  \draw[fl] ([yshift=-2mm]gw.west) -- node[lbl,below=1pt]{4 Hz} ([yshift=-2mm]rb.east);
  \draw[fl] ([yshift=2mm]gw.east) -- node[lbl,above=1pt]{action} ([yshift=2mm]nav.west);
  \draw[fl] ([yshift=-2mm]nav.west) -- node[lbl,below=1pt]{retour} ([yshift=-2mm]gw.east);
  \draw[fl] (nav) -- node[lbl]{\code{cmd\_vel}} (gz);
  \draw[fl] (gz) -- node[lbl]{capteurs, AMCL} (tf);
  \draw[fl] (tf) -- node[lbl]{pose, coûts} (gw);
  \draw[fl] (ui) -- node[lbl]{HTTP REST} (api);
\end{tikzpicture}}
\caption{Architecture de la chaîne de commande et de retour (en rouge, le nœud développé).}
\label{fig:architecture}
\end{figure}

\section{Choix d'une passerelle dédiée}
Deux approches étaient possibles pour envoyer une destination depuis le navigateur :
publier directement un message sur le topic \code{/goal\_pose} ou appeler l'action Nav2
depuis le navigateur, ou bien passer par un nœud intermédiaire. La première approche est
simple mais ne fournit aucun accusé de réception : le navigateur ne sait pas si la commande
a été reçue, refusée ou perdue. La seconde place dans le navigateur toute la logique de
vérification, qui peut alors être contournée ou dupliquée par plusieurs onglets.

La passerelle dédiée a été retenue pour les raisons suivantes :
\begin{itemize}
  \item chaque commande reçoit une réponse explicite (\emph{acceptée} ou \emph{refusée}
    avec un motif) ;
  \item les vérifications (monde, coordonnées, carte, disponibilité de Nav2) sont faites
    côté ROS, au plus près des données ;
  \item un seul but peut être actif à la fois, quel que soit le nombre de navigateurs ;
  \item l'état publié est unique et partagé : tous les clients voient la même information.
\end{itemize}

\section{Protocole de communication}

\subsection{Service de soumission}
La soumission utilise une interface de service dédiée, définie dans le paquet
\code{navigation\_msgs} :
\begin{lstlisting}
# Coordinates in the dashboard's configured SDF world (metres, radians).
string request_id
string world_id
float64 x
float64 y
float64 yaw
---
bool accepted
string message
\end{lstlisting}
Le champ \code{request\_id} est un identifiant aléatoire de 32 caractères hexadécimaux
généré par le navigateur. Il permet à la passerelle de reconnaître une demande déjà reçue
et de ne pas l'exécuter deux fois. Le champ \code{world\_id} identifie le monde affiché dans
le dashboard ; la passerelle refuse la demande s'il ne correspond pas au monde simulé.
L'annulation utilise le service standard \code{std\_srvs/srv/Trigger}.

Un acquittement positif signifie uniquement \og demande reçue et transmise à Nav2 \fg{}. Le
résultat réel de la navigation n'est connu qu'à travers l'état publié.

\subsection{État de navigation}
La passerelle publie à 4~Hz sur \code{/dashboard/navigation/state} un message
\code{std\_msgs/String} contenant un objet JSON versionné (tableau~\ref{tab:etat}). Ce
choix évite de créer un type de message supplémentaire côté navigateur tout en gardant un
format validé strictement par le frontend.

\begin{table}[htbp]
\centering\small
\begin{tabularx}{\linewidth}{@{}lX@{}}
\toprule
\textbf{Champ} & \textbf{Contenu} \\
\midrule
\code{version} & Version du format (1) ; toute autre valeur est rejetée par le frontend. \\
\code{world\_id} & Monde configuré sur la passerelle. \\
\code{request\_id} & Identifiant de la dernière demande acceptée. \\
\code{phase}, \code{detail} & Phase de la demande et message lisible. \\
\code{ready}, \code{unavailable\_reason} & Disponibilité pour une nouvelle destination et motif sinon. \\
\code{pose} & Position $(x, y)$ et cap du robot dans le repère du monde, ou \code{null}. \\
\code{target} & Destination en cours, dans le repère du monde. \\
\code{distance\_remaining} & Distance restante fournie par le retour de Nav2. \\
\bottomrule
\end{tabularx}
\caption{Contenu de l'état de navigation publié par la passerelle.}
\label{tab:etat}
\end{table}

\subsection{Séquence d'une navigation}
La figure~\ref{fig:sequence} montre les échanges lors d'une navigation réussie.

\begin{figure}[htbp]
\centering
\begin{tikzpicture}[y=-0.62cm]
  \foreach \x/\n in {0/Navigateur,3.6/Passerelle,7.2/Nav2,10.4/Gazebo}{
    \node[bloc,fill=gray!8,minimum width=2.4cm] at (\x,0) {\n};
    \draw[gray,dashed] (\x,0.8) -- (\x,12.2);}
  \draw[fl] (0,1.5) -- node[lbl,above]{submit(request\_id, x, y, cap)} (3.6,1.5);
  \node[font=\scriptsize\itshape,align=left,anchor=east] at (3.5,2.4) {vérifications};
  \draw[fl] (3.6,3.2) -- node[lbl,above]{send\_goal} (7.2,3.2);
  \draw[fl,dashed] (3.6,3.9) -- node[lbl,above]{accepted = true} (0,3.9);
  \draw[fl,dashed] (7.2,4.8) -- node[lbl,above]{but accepté} (3.6,4.8);
  \draw[fl] (7.2,5.6) -- node[lbl,above]{cmd\_vel} (10.4,5.6);
  \draw[fl,dashed] (7.2,6.6) -- node[lbl,above]{feedback (distance)} (3.6,6.6);
  \draw[fl,dashed] (3.6,7.4) -- node[lbl,above]{état : navigating} (0,7.4);
  \draw[fl,dashed] (7.2,9.2) -- node[lbl,above]{résultat : SUCCEEDED} (3.6,9.2);
  \draw[fl,dashed] (3.6,10.2) -- node[lbl,above]{état : succeeded} (0,10.2);
  \node[font=\scriptsize\itshape,anchor=west] at (3.7,11.3) {état publié à 4 Hz en continu};
\end{tikzpicture}
\caption{Séquence d'une navigation réussie.}
\label{fig:sequence}
\end{figure}

\section{Machine d'états d'une demande}
Chaque demande suit la machine d'états de la figure~\ref{fig:etats}. Les phases
\code{sending}, \code{navigating} et \code{canceling} sont dites \emph{actives} : pendant
ces phases, toute nouvelle destination est refusée. Une annulation demandée pendant la
phase \code{sending}, avant l'acceptation par Nav2, est mémorisée et transmise dès que le
but est accepté. Si Nav2 refuse l'annulation, la phase revient à \code{navigating} avec un
message explicite, afin de ne jamais laisser croire que le robot s'est arrêté.

\begin{figure}[htbp]
\centering
\begin{tikzpicture}[st/.style={bloc,rounded corners=8pt,minimum width=2.2cm},
  node distance=1.1cm and 1.5cm]
  \node[st] (idle) {idle};
  \node[st,fill=yellow!15,right=of idle] (send) {sending};
  \node[st,fill=yellow!15,right=of send] (nav) {navigating};
  \node[st,fill=yellow!15,below=1.4cm of nav] (can) {canceling};
  \node[st,fill=green!15,right=1.3cm of nav] (ok) {succeeded};
  \node[st,fill=red!10,above=1.1cm of nav] (fail) {failed};
  \node[st,fill=red!10,above=1.1cm of send] (rej) {rejected};
  \node[st,fill=gray!15,right=1.3cm of can] (cand) {canceled};
  \draw[fl] (idle) -- node[lbl,above]{submit} (send);
  \draw[fl] (send) -- node[lbl,above]{accepté} (nav);
  \draw[fl] (send) -- node[lbl,left]{refusé} (rej);
  \draw[fl] (nav) -- node[lbl,above]{succès} (ok);
  \draw[fl] (nav) -- node[lbl,left]{échec} (fail);
  \draw[fl] (nav) -- node[lbl,left]{cancel} (can);
  \draw[fl] (send) |- node[lbl,pos=0.3,left]{cancel} (can);
  \draw[fl] (can) -- node[lbl,above]{annulé} (cand);
\end{tikzpicture}
\caption{Phases d'une demande de navigation (en jaune, les phases actives).}
\label{fig:etats}
\end{figure}

\section{Gestion des repères}
Le dashboard dessine le monde à partir du fichier SDF de Gazebo ; ses coordonnées sont
donc celles du monde simulé. Nav2, lui, travaille dans le repère \code{map} défini par la
carte chargée. Ces deux repères ne coïncident pas nécessairement : une carte construite par
SLAM a pour origine la position du robot au début de l'enregistrement.

La passerelle est configurée par la pose de l'origine du repère \code{map} dans le monde,
notée $(t_x, t_y, \theta)$ (paramètres \code{world\_offset\_x}, \code{world\_offset\_y},
\code{world\_offset\_yaw}). La conversion d'une pose de la carte vers le monde est :
\begin{equation}
\begin{pmatrix} x_w \\ y_w \end{pmatrix} =
\begin{pmatrix} \cos\theta & -\sin\theta \\ \sin\theta & \cos\theta \end{pmatrix}
\begin{pmatrix} x_m \\ y_m \end{pmatrix} +
\begin{pmatrix} t_x \\ t_y \end{pmatrix},
\qquad \psi_w = \psi_m + \theta,
\label{eq:map2world}
\end{equation}
où $(x_m, y_m, \psi_m)$ est la pose dans \code{map} et $(x_w, y_w, \psi_w)$ la pose dans
le monde. Une destination choisie dans le dashboard est convertie par la transformation
inverse avant d'être envoyée à Nav2 :
\begin{equation}
\begin{pmatrix} x_m \\ y_m \end{pmatrix} =
\begin{pmatrix} \cos\theta & \sin\theta \\ -\sin\theta & \cos\theta \end{pmatrix}
\begin{pmatrix} x_w - t_x \\ y_w - t_y \end{pmatrix},
\qquad \psi_m = \psi_w - \theta.
\label{eq:world2map}
\end{equation}
Le cap est manipulé en radians côté ROS et saisi en degrés dans l'interface. La position
du robot est toujours lue par TF dans le repère \code{map}, jamais à partir de l'odométrie
seule, qui dérive.

\section{Vérification des destinations}
Avant d'envoyer un but à Nav2, la passerelle consulte la costmap globale. La destination,
convertie dans \code{map}, est ramenée dans le repère de la grille (origine et rotation de
la grille), puis la cellule est obtenue par
\begin{equation}
  i = \left\lfloor \frac{g_x}{r} \right\rfloor, \qquad
  j = \left\lfloor \frac{g_y}{r} \right\rfloor,
\end{equation}
où $(g_x, g_y)$ sont les coordonnées dans la grille et $r$ sa résolution. L'arrondi par
défaut (et non la troncature) garantit qu'un point situé juste avant un bord négatif est
bien considéré comme hors de la grille. La destination est refusée si la cellule est hors
carte, inconnue ($-1$) ou si son coût atteint le seuil d'occupation de 65 (sur une échelle
de 0 à 100). La carte approximative dessinée par le dashboard n'est jamais utilisée pour ce
contrôle. Une cellule libre peut toutefois être inaccessible : dans ce cas, c'est Nav2 qui
renvoie un échec, et l'interface l'affiche sans le transformer.

\section{Règles de sûreté et de fraîcheur}
Le tableau~\ref{tab:regles} résume les règles appliquées de part et d'autre.

\begin{table}[htbp]
\centering\small
\begin{tabularx}{\linewidth}{@{}lX@{}}
\toprule
\textbf{Côté} & \textbf{Règle} \\
\midrule
Passerelle & Refus si le temps simulé n'est pas utilisé (\code{use\_sim\_time} faux). \\
Passerelle & Refus si Nav2 est absent, si la transformation TF a plus de 2~s ou si la costmap a plus de 5~s. \\
Passerelle & Un seul but actif ; un \code{request\_id} déjà reçu n'est jamais réexécuté. \\
Passerelle & Horloge de publication basée sur le temps simulé : l'état n'est plus publié si Gazebo est en pause. \\
Dashboard & État considéré périmé après 2,5~s : position retirée de la carte et commandes bloquées. \\
Dashboard & Commandes interdites en mode démonstration ou si le monde affiché diffère du monde simulé. \\
Dashboard & Délai d'acquittement de 6~s ; en cas d'expiration, aucun renvoi automatique et message invitant à consulter l'état. \\
\bottomrule
\end{tabularx}
\caption{Règles de sûreté et de fraîcheur des données.}
\label{tab:regles}
\end{table}

\section*{Conclusion}
La conception repose sur une passerelle ROS~2 unique, un protocole acquitté et un état
publié en continu, ainsi que sur une conversion explicite entre le repère de la carte et
celui du monde. Le chapitre suivant présente la réalisation de ces éléments.
